\chapter{Fourier Series}

\section{Jensen Inequality}

\begin{definitionenv}
    Let $E$ be a vector space over $\RR$, $S \subseteq E$ be a subset.
    We call a \textbf{convex combination} of elements of $S$ any element $x \in E$ which can be written as
    $$x = \lambda_1 x_1 + \cdots + \lambda_n x_n,$$
    where $n$ is a positive integer, $\lambda_1,\cdots,\lambda_n$ are non-negative real numbers such that
    $$\lambda_1 + \cdots + \lambda_n = 1$$
    and $(x_1,\cdots,x_n) \in S^n$.

    \quad We denote by $\Conv(S)$ the set of all convex combinations of elements of $S$ and call it the \textbf{convex hull} of $S$.
    If $\Conv(S) = S$, then we say that $S$ is a convex subset of $E$.

    \quad Let $U$ be a convex subset of $E$, $f : U \longrightarrow \RR$ be a mapping.
    If the epigraph of $f$, i.e.
    $$\Gamma_{+} = \{(x,t) \in U \times \RR \mid f(x) \le t\}$$
    is a convex subset of $E \times \RR$, then we say that $f$ is \textbf{convex}.
    If the hypograph of $f$, i.e.
    $$\Gamma_{-} = \{(x,t) \in U \times \RR \mid f(x) \ge t\}$$
    is a convex subset of $E \times \RR$, then we say that $f$ is \textbf{concave}.
\end{definitionenv}

\begin{propositionenv}
    Let $E$ be a vector space over $\RR$, $U \subseteq E$ be a subset.
    Then $U$ is convex if and only if for any $t \in [0,1]$, any $(x,y) \in U^2$,
    $$tx + (1-t)y \in U.$$
\end{propositionenv}
\begin{proofenv}
    The necessary is clear. For the sufficiency, we reason by induction with respect to the number of elements in the convex combination denoted as $n$.
    The cases $n=1$ and $n=2$ are clear.
    In the following, we prove the passage from $n$ to $n+1$ elements in the convex combination.
    Let $x_1,\cdots,x_{n+1} \in U$ and $\lambda_1,\cdots,\lambda_{n+1} \in \RR_{\ge 0}$, $\lambda_1 + \cdots + \lambda_{n+1} = 1$.
    If $\lambda_{n+1}=1$, the assertion is clear. Otherwise $\lambda_1+\cdots+\lambda_n>0$.
    For $i \in \{1,\cdots,n\}$, put
    $\mu_i = \frac{\lambda_i}{\lambda_1 + \cdots + \lambda_n}$.
    Then $\mu_1 + \cdots + \mu_n = 1$.
    By the induction hypothesis, we obtain
    $$y \coloneq \mu_1 x_1 + \cdots + \mu_n x_n \in U.$$
    Let $t = \lambda_1 + \cdots + \lambda_n$.
    One has $1-t = \lambda_{n+1}$.
    By the case $n=2$, $ty + (1-t)x_{n+1} = \lambda_{1}x_1 + \cdots + \lambda_{n+1}x_{n+1} \in U$ as required.
\end{proofenv}

\begin{theoremenv}[Jensen's Inequality]
    Let $E$ be a vector space over $\RR$, $U \subseteq E$ be a convex subset,
    $f : U \longrightarrow \RR$ be a mapping, then the following statements are equivalent:
    \begin{enumerate}[label=(\arabic*)]
        \item $f$ is convex.
        \item For any $n \in \NN_{\ge 1}$, any $(x_1,\cdots,x_n) \in U^n$, any $(\lambda_1,\cdots,\lambda_n) \in \RR_{\ge 0}^n$ such that $\lambda_1 + \cdots + \lambda_n = 1$, one has
            $$f(\lambda_1 x_1 + \cdots + \lambda_n x_n) \le \lambda_1 f(x_1) + \cdots + \lambda_n f(x_n).$$
        \item  For any $(x,y) \in U^2$, any $t \in [0,1]$, one has
            $$f(tx + (1-t)y) \le t f(x) + (1-t)f(y).$$
    \end{enumerate}
\end{theoremenv}
\begin{proofenv}
    The inequality
    $$f(\lambda_1 x_1 + \cdots + \lambda_n x_n) \le \lambda_1 f(x_1) + \cdots + \lambda_n f(x_n)$$
    can be interpreted as
    $$(\lambda_1 x_1 + \cdots + \lambda_n x_n, \lambda_1 f(x_1) + \cdots + \lambda_n f(x_n)) = \lambda_1 (x_1,f(x_1)) + \cdots + \lambda_n (x_n,f(x_n))$$
    belongs to the epigraph of $f$.
    The first two statements are therefore equivalent.
    The first and the third statements are equivalent by the similar reasoning.
\end{proofenv}

\begin{definitionenv}
    Let $E$ be a vector space over $\RR$.
    We call \textbf{affine function} on $E$ any mapping $f : E \longrightarrow \RR$ of the form
    $$f(x) = \varphi(x) + C,$$
    where $\varphi : E \longrightarrow \RR$ is $\RR$-linear and $C \in \RR$ is a constant.
    For any $(x,y) \in E \times E$, $t \in [0,1]$,
    $$\begin{aligned}
        f(tx + (1 - t)y)
        &= \varphi (tx + (1 - t)y) + C\\
        &= t \varphi(x) + (1 - t)\varphi(y) + C\\
        &= t \varphi(x) + (1 - t)\varphi(y) + C(t + (1 - t))\\
        &= t f(x) + (1 - t) f(y).
    \end{aligned}$$
    So $f$ is both convex and concave.
\end{definitionenv}

\begin{propositionenv}
    Let $E$ be a vector space over $\RR$, $U \subseteq E$ be a convex subset.
    Let $(f_i)_{i \in I}$ be a family of convex mappings from $U$ to $\RR$.
    Assume that for any $x \in U$, $\displaystyle \sup_{i \in I}f_i(x) \in \RR$, then the mapping
    $$\begin{array}{rrcl}
        f : &  U & \longrightarrow & \RR\\
        & x & \longmapsto & \displaystyle \sup_{i \in I}f_i(x)
    \end{array}$$
    is convex.
\end{propositionenv}
\begin{proofenv}
    Let $(x,y) \in U^2$, $t \in [0,1]$.
    $$\begin{aligned}
        f(tx + (1 - t)y) &= \sup_{i \in I}f_i(tx + (1 - t)y)\\
        & \le \sup_{i \in I}(t f_i(x) + (1 - t)f_i(y))\\
        & \le \sup_{i \in I} t f_i(x) + \sup_{i \in I}(1 - t)f_i(y)\\
        & = t \sup_{i \in I}f_i(x) + (1 - t) \sup_{i \in I}f_i(y)\\
        & = t f(x) + (1 - t) f(y).
    \end{aligned}$$
\end{proofenv}

\begin{propositionenv}
    Let $H$ be a Hilbert space,
    $U \subseteq H$ be a convex subset,
    $f : U \longrightarrow \RR$ be a mapping.
    Assume that $f$ is at least twice differentiable and $\DD^2 f(x)$ is semipositive (resp. seminegative) for any $x \in U$.
    Then $f$ is convex (resp. concave).
    \quad Moreover, for any $a \in U$,
    $$\forall x \in U,\quad f(x) \ge f(a) + \DD f(a)(x - a) \quad \left(\text{resp. } f(x) \le f(a) + \DD f(a)(x - a) \right).$$
\end{propositionenv}
\begin{proofenv}
    It suffices to prove the semipositive case.
    We claim that for arbitrary $a \in U$, one has
    $$\forall x \in U \quad f(x) \ge f(a) + \DD f(a)(x - a).$$
    Consider
    $$g : t \longmapsto f(a + t(x - a)).$$
    $g$ is at least twice differentiable.
    One has
    $$g'(t) = \DD f(a + t(x - a))(x - a),$$
    $$g''(t) = \DD^2 f(a + t(x - a))(x - a, x - a).$$
    By the Taylor's formula (Lagrange form), there exists $\theta \in [0,1]$ such that
    $$\begin{aligned}
        f(x) = & g(1) = g(0) + g'(0) + \frac{1}{2} g''(\theta)\\
        & = f(a) + \DD f(a)(x - a) + \frac{1}{2} \DD^2 f(a + \theta(x - a))(x - a, x - a)\\
        & \ge f(a) + \DD f(a)(x - a). 
    \end{aligned}$$
    Denote by $f_a$ the affine function defined as
    $$f_a(x) \coloneq f(a) + \DD f(a) (x - a).$$
    Then
    $$f = \sup_{a \in U} f_{a}.$$
    Hence $f$ is convex.
\end{proofenv}

\section{Hölder Inequality}
We fix a measure space $(\Omega,\mathscr{A},\nu)$, and $\mathbb{K} \in \{\RR,\CC\}$.

\begin{definitionenv}
    Let $p \in \interval[open right]{1}{\infty}$.
    If $f : \Omega \longrightarrow \mathbb{K}$ is $\mathscr{A}$-measurable, we put
    $$\|f\|_{L^{p}} = \left(\int_{\Omega} |f(\omega)|^p \nu(\dd \omega)\right)^{1/p}.$$
    We denote by $\mathcal{L}^{p}(\Omega,\mathscr{A},\nu)$ the set of $\mathscr{A}$-measurable mappings from $\Omega$ to $\mathbb{K}$ such that $\|f\|_{L^{p}} < + \infty$.

    \quad For $\lambda \in \mathbb{K}$,
    $$\|\lambda f\|_{L^{p}} = |\lambda| \|f\|_{L^{p}}$$
    provided that $\lambda f \in \mathcal{L}^{p}(\Omega,\mathscr{A},\nu)$.
\end{definitionenv}

\begin{definitionenv}
    For any $f : \Omega \longrightarrow \mathbb{K}$ measurable, put
    $$\|f\|_{L^{\infty}} = \inf_{\substack{A \in \mathscr{A}\\\nu(A) = 0}}\sup_{\omega \in \Omega \setminus A} |f(\omega)|.$$
    We denote by $\mathcal{L}^{\infty}(\Omega,\mathscr{A},\nu)$ the set of all measurable $f : \Omega \longrightarrow \mathbb{K}$ such that $\|f\|_{L^{\infty}} < +\infty$.
    Note that the set
    $$\{\omega \in \Omega \mid |f(\omega)| > \|f\|_{L^{\infty}}\} = \bigcup_{\varepsilon \in \QQ_{>0}}\{\omega \in \Omega \mid |f(\omega)| > \|f\|_{L^{\infty}} + \varepsilon\}$$
    is negligible.

    \quad It turns out that $\mathcal{L}^{\infty}(\Omega,\mathscr{A},\nu)$ forms a vector space over $\mathbb{K}$ and $\|\cdot\|_{L^{\infty}}$ defines a seminorm on $\mathcal{L}^{\infty}(\Omega,\mathscr{A},\nu).$

    \quad Indeed, if $f,g \in \mathcal{L}^{\infty}(\Omega,\mathscr{A},\nu)$ and $A \in \mathscr{A}$, $\nu(A) = 0$,
    $$\sup_{\omega \in \Omega\setminus A} |(f + g)(\omega)| \le (\sup_{\omega \in \Omega \setminus A}|f(\omega)|) + (\sup_{\omega \in \Omega \setminus A}|g(\omega)|).$$
    Taking the infimum with respect to $A \in \mathscr{A}$, we obtain
    $$\|f + g\|_{L^{\infty}} \le \|f\|_{L^{\infty}} + \|g\|_{L^{\infty}}.$$

    \quad The null space of the seminorm $\|\cdot\|_{L^{\infty}}$ is the set of all negligible functions.
    We denote by $L^{\infty}(\Omega,\mathscr{A},\nu)$ the quotient of $\mathcal{L}^{\infty}(\Omega,\mathscr{A},\nu)$ by the null space.
    The seminorm $\|\cdot\|_{L^{\infty}}$ induces a norm on $L^{\infty}(\Omega,\mathscr{A},\nu)$.
    By abuse of notation, this norm is denoted by $\|\cdot\|_{L^{\infty}}$.
\end{definitionenv}

\begin{theoremenv}[Hölder Inequality]
    Let $n$ be a positive integer, $p_1,\cdots,p_n \in \RR_{\ge 1} \cup \{\infty\}$ such that
    $$\frac{1}{p_1} + \cdots + \frac{1}{p_n} = 1,$$
    where $\frac{1}{\infty} \coloneq 0$ by convention.
    Let $f_1,\cdots,f_n : \Omega \longrightarrow \mathbb{K}$ be measurable.
    One has
    $$\|\prod_{i=1}^{n}f_i\|_{L^{1}} \le \prod_{i = 1}^n \|f_i\|_{L^{p_i}},$$
    provided that $\dis \prod_{i = 1}^n \|f_i\|_{L^{p_i}}$ is well defined.
\end{theoremenv}
\begin{proofenv}
    We may assume that $\|f_{j}\|_{L^{p_j}} > 0$ for any $j \in \{1,\cdots,n\}$, since otherwise
    $$\{\omega \in \Omega \mid f_{j} (\omega) \neq 0\}$$
    is of measure $0$ and hence $f_1 \cdot \cdots \cdot f_n$ is also negligible, which leads to
    $$\|f_1 \cdot \cdots \cdot f_n\|_{L^{1}} = 0.$$
    We reason by induction on $n$.
    The case $n=1$ is clear.
    Assume that the statement holds for the case $n-1$. If at least one of $p_1, \cdots, p_n$ is $\infty$, let $p_n = \infty$, then
    $$\frac{1}{p_1} + \cdots + \frac{1}{p_{n-1}} = 1.$$
    So the thesis follows from the induction hypothesis.
    In the remaining part, $p_1,\cdots,p_n \in \RR_{\ge 1}$.
    For any $x_1,\cdots,x_n \in \RR_{>0}$,
    $$\prod_{i = 1}^{n} x_i = \exp \left(\sum_{j = 1}^{n} \frac{\ln(x_j^{p_j})}{p_j}\right).$$
    Since $\exp$ is convex and $\frac{1}{p_1} + \cdots + \frac{1}{p_n} = 1$, by Jensen's inequality,
    $$\prod_{i = 1}^{n} x_i \le \sum_{j = 1}^{n} \frac{1}{p_j} \exp\left(\ln(x_j^{p_j})\right) = \sum_{j=1}^{n} \frac{x_j^{p_j}}{p_j}.$$
    It follows that
    $$\frac{|\prod_{i=1}^{n} f_i|}{\prod_{i=1}^{n}\|f_i\|_{L^{p_i}}} \le \sum_{j=1}^{n} \frac{|f_j|^{p_j}}{p_j \cdot \|f_j\|^{p_j}_{L^{p_j}}}.$$
    Taking the integral with respect to $\nu$,
    $$\frac{1}{\prod_{i=1}^{n} \|f_i\|_{L^{p_i}}} \int_{\Omega} |\prod_{i=1}^{n} f_i| (\omega) \dd \nu(\omega) \le \sum_{i=1}^{n}\frac{1}{p_i} = 1.$$
\end{proofenv}

\begin{theoremenv}[Minkowski Inequality]\label{thm:8.2.4}
    Let $p \in \RR_{\ge 1}$ and $f,g \in \mathcal{L}^{p}(\Omega,\mathscr{A},\nu)$, then $f + g \in \mathcal{L}^{p}(\Omega,\mathscr{A},\nu)$,
    $$\|f + g\|_{L^{p}} \le \|f\|_{L^{p}} + \|g\|_{L^{p}}.$$
    In particular, $\mathcal{L}^{p}(\Omega,\mathscr{A},\nu)$ is a vector space over $\mathbb{K}$ and $\|\cdot\|_{L^{p}}$ is a seminorm on $\mathcal{L}^{p}(\Omega,\mathscr{A},\nu).$
\end{theoremenv}
\begin{proofenv}
    If $p = 1$, the statement follows from $|f + g| \le |f| + |g|$.
    Suppose that $p > 1$, we put $q \coloneq (1 - p^{-1})^{-1}$ so that $p^{-1} + q^{-1} = 1$.
    Without loss of generality, $f,g$ only take nonnegative values.
    Apply Hölder's inequality to $f$ and $(f + g)^{p-1}$,
    $$\int_{\Omega} f(f + g)^{p-1} \dd \nu \le \|f\|_{L^p} \left(\int_{\Omega} (f + g)^{p} \dd \nu\right)^{\frac{p-1}{p}}.$$
    Similarly,
    $$\int_{\Omega} g(f + g)^{p-1} \dd \nu \le \|g\|_{L^p} \left(\int_{\Omega} (f + g)^{p} \dd \nu\right)^{\frac{p-1}{p}}.$$
    Summing the two inequalities, we obtain
    $$\|f+g\|^p_{L^p} = \int_{\Omega} (f + g)^{p} \dd \nu \le (\|f\|_{L^p} + \|g\|_{L^p}) \|f+g\|_{L^p}^{p-1}.$$
\end{proofenv}

\begin{remark}
    Note that the null space of $\mathcal{L}^p(\Omega,\mathscr{A},\nu)$ with respect to the seminorm $\|\cdot\|_{L^p}$ coincides with the vector space of negligible functions.
    We denote by $L^p(\Omega,\mathscr{A},\nu)$ the quotient of $\mathcal{L}^p(\Omega,\mathscr{A},\nu)$ by the negligible functions.
    The seminorm $\|\cdot\|_{L^p}$ induces a norm on $L^p(\Omega,\mathscr{A},\nu)$.
\end{remark}

\begin{lemmaenv}\label{lem:8.2.6}
    Let $p \in \RR_{\ge 1} \cup \{\infty\}$ and $(u_n)_{n \in \NN}$ be a sequence of elements in $\mathcal{L}^{p}(\Omega,\mathscr{A},\nu)$.
    Suppose that
    $$C \coloneq \sum_{n \in \NN}\|u_n\|_{L^p} < + \infty.$$
    Then there exists $A \in \mathscr{A}$ such that $\nu(A) = 0$ and $(u_n(\omega))_{n \in \NN}$ converges absolutely for $\omega \in \Omega \setminus A$.
    Moreover, if we denote by
    $$\begin{array}{rrcl}
        u : & \Omega & \longrightarrow & \RR\\
        & \omega & \longmapsto & \begin{cases}
            \dis \sum_{n \in \NN} u_n(\omega) & \text{if } \omega \in \Omega \setminus A\\
            0 & \text{if } \omega \in A
        \end{cases},
    \end{array}$$
    then $u \in \mathcal{L}^p(\Omega,\mathscr{A},\nu)$ and
    $$\lim_{n \rightarrow +\infty}\|u - \sum_{k=0}^{n}u_k\|_{L^p} = 0.$$
\end{lemmaenv}
\begin{proofenv}
    For any $n \in \NN$,
    $$v_n = \sum_{k=0}^{n}|u_k|, \quad v = \sum_{n \in \NN}|u_n| = \lim_{n \rightarrow \infty} v_n.$$
    The sequence $(v_n)_{n \in \NN}$ is increasing.
    For $p < \infty$, by \ref{thm:8.2.4},
    $$\|v_n\|_{L^p} \le \sum_{k=0}^{n} \|u_k\|_{L^p} \le C.$$
    For $p = \infty$,
    $$\{\omega \in \Omega \mid v_n(\omega) > C\} \subseteq \{\omega \in \Omega \mid v(\omega) > C\},$$
    $$\{\omega \in \Omega \mid v(\omega) = +\infty\} \subseteq \bigcup_{n \in \NN} \{\omega \in \Omega \mid v_{n}(\omega) > C\}.$$
    Hence $A \coloneq \{\omega \in \Omega \mid v(\omega) = + \infty\}$ is of measure $0$.
    We can deduce that for any $\omega \in \Omega \setminus A$, $\sum_{n \in \NN} u_n (\omega)$ converges absolutely.
    Put
    $$u(\omega) = \begin{cases}
        \sum_{n \in \NN} u_n(\omega), & \omega \in \Omega \setminus A\\
        0, & \omega \in A
    \end{cases}.$$
    Since $|u|\le v$ for any $\omega \in \Omega$, we obtain
    $$u \in \mathcal{L}^p(\Omega,\mathscr{A},\nu).$$
    For $p \neq \infty$,
    $$\|u\|_{L^p} \le \left(\int_{\Omega} v(\omega)^p \dd \nu(\omega)\right)^{\frac{1}{p}} \le \sum_{n \in \NN} \|u_n\|_{L^p}.$$
    For $p = \infty$,
    $$\|u\|_{L^\infty} \le \|v\|_{L^\infty} \le \sum_{n \in \NN} \|u_n\|_{L^\infty}.$$
    Similarly,
    $$\|u - \sum_{k=0}^{n}u_k\|_{L^p} \le \sum_{k=n+1}^{\infty}\|u_k\|_{L^p},$$
    so
    $$\lim_{n \rightarrow \infty} \|u - \sum_{k=0}^{n} u_k\|_{L^p} = 0.$$
\end{proofenv}

\begin{propositionenv}\label{prop:8.2.7}
    Let $p \in \RR_{\ge 1}$.
    The seminormed vector space $\mathcal{L}^p(\Omega,\mathscr{A},\nu)$ is complete.
\end{propositionenv}
\begin{proofenv}
    Let $(f_n)_{n \in \NN}$ be a Cauchy sequence in $\mathcal{L}^p(\Omega, \mathscr{A}, \nu)$.
    We will extract a subsequence $(f_{n_{k}})_{k \in \NN}$ in the following way:
    if $n_0,\cdots,n_k$ are chosen, then $n_{k+1}$ is chosen so that $n_{k+1} > n_k$ and
    $$\forall l,m \in \NN_{\ge n_{k+1}}, \; \|f_l - f_m\| \le 2^{-(k+1)}.$$
    Thus
    $$\|f_{n_0}\|_{L^p} + \sum_{k \in \NN} \|f_{n_{k+1}} - f_{n_k}\|_{L^p} \le \|f_{n_0}\|_{L^p} + \sum_{k \in \NN}2^{-k} = \|f_{n_0}\|_{L^p} + 2 < + \infty.$$
    By lemma \ref{lem:8.2.6}, $(f_{n_{k}})_{k \in \NN}$ converges in $\mathcal{L}^p(\Omega,\mathscr{A},\nu)$.
\end{proofenv}

\begin{corollaryenv}
    The $\mathbb{K}$-vector space $L^2(\Omega,\mathscr{A},\nu)$ with the inner product
    $$\langle f,g \rangle_{L^2} = \int_{\Omega} \bar{f}(\omega) g(\omega) \dd \nu(\omega)$$
    forms a Hilbert space.
\end{corollaryenv}
\begin{remark}
    From the proofs of lemma \ref{lem:8.2.6} and proposition \ref{prop:8.2.7},
    we observe that if $(f_n)_{n \in \NN}$ is a convergent sequence in $\mathcal{L}^p(\Omega,\mathscr{A},\nu)$,
    then it has a subsequence $(f_{n_k})_{k \in \NN}$ such that the series
    $$\sum_{k \in \NN_{\ge 1}}(f_{n_{k}} - f_{n_{k-1}})$$
    converges absolutely $\nu$-almost everywhere to the limit of $(f_n)_{n \in \NN}$ in $\mathcal{L}^p(\Omega,\mathscr{A},\nu)$.
\end{remark}

\begin{remark}
    Assume that $\nu(\Omega) < +\infty$.
    For any $p,q \in \RR_{\ge 1} \cup \{\infty\}$ such that $p < q$, one has
    $$\mathcal{L}^q(\Omega,\mathscr{A},\nu) \subseteq \mathcal{L}^p(\Omega,\mathscr{A},\nu).$$
    Indeed, if $q \in \RR_{\ge 1}$, we have $\mathbbm{1}_{\Omega} \in \mathcal{L}^{q/(q - p)}(\Omega,\mathscr{A},\nu)$.
    From Hölder (for $|f|^p$ and $\mathbbm{1}_{\Omega}$)
    $$\int_{\Omega} |f|^p \dd \nu \le \left(\int_{\Omega} |f|^q \dd \nu\right)^{p/q} \cdot \left(\int_{\Omega} \mathbbm{1}_{\Omega}^{q/(q-p)}\right)^{(q-p)/q}.$$
    Hence $\|f\|^{p}_{L^p} \le \|f\|^p_{L^q} \cdot \nu(\Omega)$.
    If $q = \infty$, for any $f \in \mathcal{L}^\infty(\Omega,\mathscr{A},\nu)$,
    $$\int_{\Omega} |f|^p \dd \nu \le \|f\|^p_{L^\infty} \nu(\Omega) < + \infty.$$
\end{remark}


\section{Haar Measure on a Torus}
Let $\mathbb{K} \in \{\RR,\CC\}$.

\begin{definitionenv}
    Let $(V,\langle \cdot,\cdot \rangle)$ be a Hilbert space over $\mathbb{K}$.
    We say that an orthogonal family $(e_i)_{i \in I}$ (resp. orthonormal family) is an orthogonal basis (resp. orthonormal basis) of $V$ if
    $$\Vect_{\mathbb{K}}(\{e_i \mid i \in I\})$$
    is dense in $V$.
\end{definitionenv}

\begin{propositionenv}
    Let $(V,\langle \cdot,\cdot \rangle)$ be a Hilbert space over $\mathbb{K}$ and $(e_i)_{i \in I}$ be an orthogonal family.
    Then $(e_i)_{i \in I}$ is an orthogonal basis if and only if
    $$\{e_i \mid i \in I\}^\perp = \{0\}.$$
\end{propositionenv}
\begin{proofenv}
    Without loss of generality, $\{e_i\}_{i \in I}$ is an orthonormal family.
    Let
    $$W = \overline{\Vect_{\mathbb{K}}(\{e_i \mid i \in I\})}.$$

    ``$\Rightarrow$'': Suppose that $W = V$.
    If $x \in V$ is such that $\langle e_i,x \rangle = 0$ for any $i \in I$, by proposition \ref{prop:3.2.7}, we obtain that $x = p_{W}(x) = 0$.

    ``$\Leftarrow$'': Suppose that $\{e_i \mid i \in I\}^{\perp} = \{0\}$.
    By proposition \ref{prop:3.2.7}, for any $x \in V$,
    $$x - p_W(x) \in W^\perp \subseteq \{e_i \mid i \in I\}^\perp.$$
    Hence $x - p_{W}(x) = 0$.
\end{proofenv}

In the rest of the section, we fix a finite-dimensional vector space $E$ over $\RR$ equipped with Haar measure $\mu$.
Let $(e_j)_{j=1}^n$ be a basis of $E$ and $(e^\vee_{j})_{j=1}^{n}$ be the dual basis.
We let
$$D \coloneq \{a_1 e_1 + \cdots + a_n e_n \mid (a_1,\cdots,a_n) \in \interval[open right]{0}{1}^n\},$$
$$\mathrm{Vol}(D) \coloneq \mu(D).$$
Let $M$ be a sub-$\ZZ$-module of $E$ generated by $(e_j)_{j=1}^n$ (i.e. $E = \Span_{\RR}(M)$), $N$ be a sub-$\ZZ$-module of $E^\vee$ generated by $(e^\vee_j)_{j=1}^{n}$ (i.e. $N = \{\varphi \in E^\vee \mid \forall x \in M, \varphi(x) \in \ZZ\}$).

For any $\alpha \in N$,
$$\begin{array}{rrcl}
    \varphi_\alpha : & E & \longrightarrow & \CC\\
    & x & \longmapsto & \exp (2 \pi \mathrm{i} \alpha(x)).
\end{array}$$
It is a group morphism from $(E,+)$ to the unit circle in $\CC$.
Note that $\alpha(m) \in \ZZ$ for any $m \in M$.
Hence $\varphi_{\alpha} (m) = 1$.
This means that $\varphi_{\alpha}$ is $M$-periodic, namely, for any $x \in E$ and $m \in M$,
$$\varphi_{\alpha}(x + m) = \varphi_{\alpha}(x).$$
In particular, $\varphi_{\alpha}$ induces a morphism of groups from $E/M$ to the unit circle in $\CC$ (denoted $\varphi_{\alpha}$ by abuse of notation).
For $\alpha, \beta \in N$,
$$\varphi_{\alpha + \beta} = \varphi_{\alpha} \varphi_{\beta}.$$
Denote by $L^2(D)$ the $\CC$ vector space of square-integrable, Borel measurable functions on $D$.
We equip $L^2(D)$ with the following inner product
$$\forall f,g \in L^2(D) \quad \langle f,g \rangle \coloneq \frac{1}{\mathrm{Vol}(D)}\int_{D} \bar{f}(x) g(x) \Vol_D(\dd x).$$
Note that this inner product is independent of the choice of Haar measure $\mu$.

\begin{propositionenv}
    The family $(\varphi_{\alpha})_{\alpha \in N}$ forms an orthonormal family in the Hilbert space $L^2(D)$.
\end{propositionenv}
\begin{proofenv}
    $E \cong \prod_{j=1}^{n} \RR e_j$.
    Suppose that the volume of $D$ is the product of the usual Haar measures on $\RR e_{j}$ (i.e. measure of $[0,1]e_j$ which is $1$).
    For any $\alpha \in N$,
    $$\langle \varphi_\alpha, \varphi_\alpha \rangle = \frac{1}{\Vol(D)} \int_D 1 \Vol_D(\dd x) = 1.$$
    For $\alpha \neq \beta$,
    $$\langle \varphi_\alpha, \varphi_\beta \rangle = \frac{1}{\Vol(D)} \int_D \exp\left(2 \pi \mathrm{i} (\beta - \alpha)(x)\right) \Vol_{D}(\dd x).$$
    Suppose that $\beta - \alpha = a_1 e^\vee_1 + \cdots + a_n e^\vee_n$, $(a_1,\cdots,a_n) \in \ZZ^n$.
    $$\int_{D} \exp(2 \pi \mathrm{i} (\beta - \alpha)(x))\Vol_D(\dd x) = \prod_{j=1}^{n} \int_{0}^{1} \exp\left(2 \pi \mathrm{i}a_{j} t\right) \dd t.$$
    If $a_j \neq 0$,
    $$\int_{0}^{1} \exp(2 \pi \mathrm{i} a_j t) \dd t = \left[\frac{\exp(2 \pi \mathrm{i} a_j t)}{2 \pi \mathrm{i} a_j}\right]_{t = 0}^{t = 1} = 0.$$
\end{proofenv}

\begin{propositionenv}
    Let $\mathbb{T}_M = E/M$, $\pi : E \longrightarrow \mathbb{T}_M$, then $\pi|_D$ is an isomorphism of measurable spaces.
    Namely, if we equip $\mathbb{T}_M$ with a Haar measure such that $\eta(\mathbb{T}_M) = 1$, then for any integrable $h : \mathbb{T}_M \longrightarrow \CC$,
    $$\int_{\mathbb{T}_M} h(t) \eta(\dd t) = \frac{1}{\Vol(D)} \int_{D} h(\pi(x)) \Vol_{D}(\dd x).$$
\end{propositionenv}
\begin{proofenv}
    $\pi|_{D}$ is a bijection and is measurable since it is continuous.
    Since the Borel $\sigma$-algebra on $D$ is generated by any topological basis, it suffices to check that for any $U \subseteq D$ of the form
    $$U = \sum_{i=1}^{d} U_i e_i,$$
    where $U_i \subseteq \interval[open right]{0}{1}$ is open.
    Let $p : \RR \longrightarrow \RR/\ZZ$ be the projection.
    Then $p(\interval[open]{0}{1} \cap U_i)$ is open in $\RR/\ZZ$,
    $p(\{0\} \cap U_i)$ is closed in $\RR/\ZZ$.
    So $p(U_i)$ is measurable.
    Therefore, $\pi(U)$ is measurable.

    \quad To check the equality, it suffices to prove that
    $$\tilde{\eta} : \mathbb{T}_M \supseteq A \longmapsto \frac{\Vol(\pi^{-1}(A) \cap D)}{\Vol(D)}$$
    is a Haar measure on $\mathbb{T}_M$.

    \quad Let $t \in \mathbb{T}_M$,
    $$\tilde{\eta}(A + t) = \frac{\Vol(\pi^{-1}(A + t) \cap D)}{\Vol(D)}.$$
    Let $x \in D$ such that $\pi(x) = t$,
    $$\pi^{-1}(A + t) \cap D= (\pi^{-1}(A) + x) \cap D,$$
    $$\pi^{-1}(A) = \bigcup_{m \in M} \left((\pi^{-1}(A) \cap D) - m\right).$$
    So,
    $$\begin{aligned}
        \pi^{-1}(A + t) \cap D = & \bigcup_{m \in M} \left((\pi^{-1}(A) \cap D) + (x - m)\right) \cap D\\
        = & \bigcup_{m \in M} \left[\left((\pi^{-1}(A) \cap D) + x\right) \cap (D + m) - m\right].
    \end{aligned}.$$
    $$\begin{aligned}
        \Vol(\pi^{-1}(A + t) \cap D) = & \sum_{m \in M} \Vol\left[\left((\pi^{-1}(A) \cap D) + x\right) \cap (D + m)\right]\\
        = & \Vol\left((\pi^{-1}(A) \cap D) + x\right) = \Vol(\pi^{-1}(A) \cap D).
    \end{aligned}$$
\end{proofenv}


\section{Convolution}

Let $(G,+)$ be a locally compact Hausdorff group such that $G$ is $\sigma$-compact (that is, $G$ is a union of countably many compact subsets).
Assume that there exists a (non-zero) Radon measure $\eta$ on $G$ that is both left and right Haar measure.

\begin{remark}
    \quad
    \begin{enumerate}[label = (\roman*)]
        \item If $(G,+)$ is commutative, then $\eta$ always exists.
        \item If $h: G \longrightarrow \RR_{\ge 0}$ is measurable
            $$\int_{G} h(x) \eta(\dd x) = \int_{G} h(- x) \eta(\dd x).$$
            This is because: if $\eta$ is a left measure, then $\tilde{\eta} : A \longmapsto \eta(-A)$ is a right Haar measure.
    \end{enumerate}
\end{remark}

\begin{definitionenv}
    Let $f$ and $g$ be measurable mappings from $G$ to $\CC$.
    For any $x \in G$, if the function
    $$(y \in G) \longmapsto f(y) g(- y + x)$$
    is integrable, we denote by $(f * g)(x)$ the integral
    $$\int_{G} f(y) g(- y + x) \eta(\dd y) = \int_{G} f(-y) g(y + x) \eta(\dd y) \overset{z = y + x}{=} \int_{G} f(x - z) g(z) \eta(\dd z).$$
\end{definitionenv}

\begin{remark}
    If $+$ is commutative, then
    $$f*g = g*f.$$
    We say that $f$ and $g$ are convolvable if the set of $x \in G$ such that $(f*g)(x)$ is NOT well defined is $\eta$-negligible.

    \quad If $f$ and $g$ are convolvable, whenever $(f*g)(x)$ is NOT well defined, by convention $(f*g)(x)$ is defined to be $0$.
\end{remark}

\begin{theoremenv}[Young Inequality]
    Let $p,q,r$ be elements of $\RR_{\ge 1} \cup \{\infty\}$ such that
    $$\frac{1}{p} + \frac{1}{q} = 1 + \frac{1}{r}.$$
    Let $f \in L^p(G)$, $g \in L^q(G)$.
    Then $f$ and $g$ are convolvable.
    Moreover,
    $$\|f * g\|_{L^r} \le \|f\|_{L^p} \cdot \|g\|_{L^q}.$$
\end{theoremenv}
\begin{proofenv}
    \quad
    \begin{enumerate}
        \item Case 1, $r = \infty$. If $q \neq \infty$, by Hölder inequality,
            $$\begin{aligned}
                & \int_{G} |f|(y) |g|(-y+x) \eta(\dd y)\\
                \le & \|f\|_{L^p} \cdot \left(\int_{G} |g(-y+x)|^q \eta(\dd y)\right)^{\frac{1}{q}}\\
                = & \|f\|_{L^p} \cdot \|g\|_{L^q}.
            \end{aligned}$$
            If $q = \infty$,
            $$\int_{G} |f(y)|\cdot|g(-y+x)| \eta(\dd y) \le \|f\|_{L^1} \cdot \|g\|_{L^\infty}.$$
            Hence $f*g \in L^\infty(G)$ and $\|f*g\|_{L^\infty(G)} \le \|f\|_{L^1(G)} \cdot \|g\|_{L^\infty(G)}.$
        \item Case 2, $p=q=r=1$.
            $$\begin{aligned}
                & \int_{G} \int_{G} |f(y)|\cdot|g(-y+x)| \eta(\dd y) \eta(\dd x)\\
                = & \int_{G} \int_{G} |f(y)|\cdot|g(-y+x)| \eta(\dd x) \eta(\dd y)\\
                = & \int_{G} |f(y)| \int_{G} |g(x)| \eta(\dd x) \eta(\dd y)\\
                = & \|f\|_{L^1} \cdot \|g\|_{L^1}.
            \end{aligned}$$
            Hence $f*g$ is convolvable and
            $$\|f * g\|_{L^1} \le \|f\|_{L^1} \cdot \|g\|_{L^1}.$$
        \item Case 3, $r \notin \{1, \infty\}$.
            Then $\infty \notin \{p,q\}$.
            Let $a,b$ in $\RR_{\ge 1} \cup \{\infty\}$ such that
            $$\frac{1}{p} + \frac{1}{a} = 1,\; \frac{1}{q} + \frac{1}{b} = 1, \; \frac{1}{a} + \frac{1}{b} + \frac{1}{r} = 1.$$
            \begin{enumerate}
                \item Suppose that $p,q,a,b$ are real numbers.
                    For any $x \in G$, by Hölder inequality,
                    $$\begin{aligned}
                        & \int_{G} |f(y)| \cdot |g(-y+x)| \eta(\dd y)\\
                        = & \int_{G} |f(y)|^{\frac{p}{b}} \cdot |g(-y+x)|^{\frac{q}{a}}
                        \cdot \left(|f(y)|^{p} \cdot |g(-y+x)|^q\right)^{\frac{1}{r}} \eta(\dd y)\\
                        \le & \|f\|_{L^p}^{\frac{p}{b}}\cdot \|g\|_{L^q}^{\frac{q}{a}}
                        \left(\int_{G} |f(y)|^p |g(-y+x)|^q \eta(\dd y)\right)^{\frac{1}{r}}.
                    \end{aligned}$$
                    Thus
                    $$\begin{aligned}
                        & \int_{G}\left(\int_{G} |f(y)| \cdot |g(-y+x)| \eta(\dd y)\right)^r \eta(\dd x)\\
                        \le & \|f\|_{L^p}^{\frac{pr}{b}} \cdot \|g\|_{L^q}^{\frac{qr}{a}}
                        \int_{G}\int_{G} |f(y)|^p |g(-y+x)|^q \eta(\dd y)\eta(\dd x)\\
                        = & \|f\|_{L^p}^{\frac{pr}{b}} \cdot \|g\|_{L^q}^{\frac{qr}{a}}
                        \int_{G} |f(y)|^p\int_{G}|g(-y+x)|^q \eta(\dd x)\eta(\dd y)\\
                        = & \|f\|_{L^p}^{\frac{pr}{b}} \cdot \|g\|_{L^q}^{\frac{qr}{a}}
                        \|f\|_{L^p}^p \cdot \|g\|_{L^q}^q\\
                        = & \|f\|_{L^p}^r \cdot \|g\|_{L^q}^r.
                    \end{aligned}$$
                \item Suppose that $p=1$. Then $q=r$.
                    For any $x \in G$, by Hölder inequality,
                    $$\begin{aligned}
                        \int_G |f(y)| |g(-y+x)| \eta(\dd y)
                        = & \int_G |f(y)|^{\frac{1}{b}}\left(|f(y)|^{\frac{1}{r}} |g(-y+x)|\right)\eta(\dd y)\\
                        \le & \|f\|_{L^1}^{\frac{1}{b}}
                        \left(\int_G |f(y)| |g(-y+x)|^r \eta(\dd y)\right)^{\frac{1}{r}}.
                    \end{aligned}$$
                    Hence
                    $$\begin{aligned}
                        & \int_G\left(\int_G |f(y)| |g(-y+x)|\eta(\dd y)\right)^r \eta(\dd x)\\
                        \le & \|f\|_{L^1}^{\frac{r}{b}}
                        \int_G\int_G |f(y)| |g(-y+x)|^r \eta(\dd y)\eta(\dd x)\\
                        = & \|f\|_{L^1}^{\frac{r}{b}} \|f\|_{L^1}\|g\|_{L^r}^r\\
                        = & \|f\|_{L^1}^r\|g\|_{L^r}^r.
                    \end{aligned}$$
                \item Suppose that $q=1$. Then $p=r$.
                    For any $x \in G$, by Hölder inequality,
                    $$\begin{aligned}
                        \int_G |f(y)| |g(-y+x)| \eta(\dd y)
                        = & \int_G |g(-y+x)|^{\frac{1}{a}}\left(|g(-y+x)|^{\frac{1}{r}} |f(y)|\right)\eta(\dd y)\\
                        \le & \|g\|_{L^1}^{\frac{1}{a}}
                        \left(\int_G |g(-y+x)| |f(y)|^r \eta(\dd y)\right)^{\frac{1}{r}}.
                    \end{aligned}$$
                    Hence
                    $$\begin{aligned}
                        & \int_G\left(\int_G |f(y)| |g(-y+x)|\eta(\dd y)\right)^r \eta(\dd x)\\
                        \le & \|g\|_{L^1}^{\frac{r}{a}}
                        \int_G\int_G |g(-y+x)| |f(y)|^r \eta(\dd y)\eta(\dd x)\\
                        = & \|g\|_{L^1}^{\frac{r}{a}} \|g\|_{L^1}\|f\|_{L^r}^r\\
                        = & \|f\|_{L^r}^r\|g\|_{L^1}^r.
                    \end{aligned}$$
            \end{enumerate}
            Therefore $f$ and $g$ are convolvable and
            $$\|f*g\|_{L^r} \le \|f\|_{L^p} \cdot \|g\|_{L^q}.$$
    \end{enumerate}
\end{proofenv}


\section{Fourier Series}
In this section, we fix a finite-dimensional vector space $E$ over $\RR$ and a sub-$\ZZ$-module $M$ of $E$ generated by a basis $\{e_i\}_{i=1}^{d}$ of $E$.
Let $N$ be the sub-$\ZZ$-module of $E^\vee$ generated by $(e^\vee_i)_{i=1}^{d}$.
Denote by $\mathbb{T}_{M}$ the quotient group $E/M$ with the unique Haar measure $\eta$, $\eta(\mathbb{T}_M) = 1$.

Moreover, for any $p \in \RR_{\ge 1} \cup \{\infty\}$, we denote by $L^p(\mathbb{T}_{M})$ the set of all measurable mappings $f:\mathbb{T}_{M} \longrightarrow \CC$ such that $\|f\|_{L^p} < +\infty$, where
$$\|f\|_{L^p} \coloneq \left(\int_{\mathbb{T}_M} |f(x)|^p \eta(\dd x)\right)^{\frac{1}{p}}.$$
If $p = \infty$, we put
$$\|f\|_{L^\infty} \coloneq \inf_{\substack{A \in \mathcal{B}(\mathbb{T}_M)\\ \eta(A) = 0}} \sup_{x \in \mathbb{T}_{M}\setminus A} |f(x)|.$$
We denote by $\pi : E \longrightarrow \mathbb{T}_{M}$ the mapping of projection.
For any $\alpha \in N$, we denote by $\varphi_{\alpha} : \mathbb{T}_M \longrightarrow \CC$ the mapping that for any $y \in E$ sends $\pi(y)$ to $\exp(2\pi \mathrm{i} \alpha(y))$.
\begin{figure}[H]
    \centering
    \begin{tikzcd}
        E & {\mathbb{T}_{M}} \\
        & {\mathbb{C}}
        \arrow["{\pi}", from=1-1, to=1-2]
        \arrow[from=1-1, to=2-2]
        \arrow["{\varphi_{\alpha}}", from=1-2, to=2-2]
    \end{tikzcd}
\end{figure}
Note that $\varphi_{\alpha} \in L^\infty(\mathbb{T}_{M})$.
We put
$$\varphi_{j} = \varphi_{e^\vee_j}.$$
Thus for any $(b_1,\cdots,b_d) \in \ZZ^d$,
$$\varphi_{b_1e^\vee_1 + \cdots + b_de^\vee_d} = \varphi_{1}^{b_1} \cdots \varphi_{d}^{b_d}.$$

\begin{definitionenv}
    Let $f \in L^{1}(\mathbb{T}_{M})$.
    For any $\alpha \in N$, we denote by $c_{\alpha}(f)$ the following complex number
    $$c_{\alpha} (f) \coloneq \int_{\mathbb{T}_{M}} f(x) \varphi_{-\alpha}(x) \eta(\dd x).$$
    The family $(c_{\alpha}(f))_{\alpha \in N}$ is called the \textbf{Fourier coefficients} of $f$.
\end{definitionenv}

\begin{propositionenv}
    Let $f,g \in L^1(\mathbb{T}_{M})$.
    One has
    $$\forall \alpha \in N, c_{\alpha} (f * g) = c_{\alpha} (f) c_{\alpha} (g).$$
\end{propositionenv}
\begin{proofenv}
    By Fubini's theorem,
    $$\begin{aligned}
        & \int_{\mathbb{T}_{M}}\exp(-2\pi \mathrm{i} \alpha(x)) \int_{\mathbb{T}_{M}} f(y) g(x - y) \eta(\dd y) \eta(\dd x)\\
        = & \int_{\mathbb{T}_{M}} f(y) \int_{\mathbb{T}_{M}} \exp(-2\pi \mathrm{i} \alpha(x)) g(x - y) \eta(\dd x) \eta(\dd y)\\
        = & \int_{\mathbb{T}_{M}} f(y) \int_{\mathbb{T}_{M}} \exp(-2\pi \mathrm{i} \alpha(x + y)) g(x) \eta(\dd x) \eta(\dd y)\\
        = & \left(\int_{\mathbb{T}_{M}} f(y) \exp(- 2\pi \mathrm{i} \alpha(y)) \eta(\dd y)\right) \left(\int_{\mathbb{T}_{M}} g(x) \exp(-2\pi \mathrm{i} \alpha(x))\eta(\dd x)\right).
    \end{aligned}$$
\end{proofenv}

\begin{definitionenv}
    We equip $E^\vee$ with the norm $\|\cdot\|$,
    $$\|\lambda_1 e^\vee_1 + \cdots + \lambda_d e^\vee_d\| = \max\{|\lambda_1|, \cdots, |\lambda_d|\}.$$
    For any $f \in L^1(\mathbb{T}_{M})$, we call \textbf{Fourier series} of the series of functions of the form
    $$S_{n} (f) \coloneq \sum_{\alpha \in N, \|\alpha\| \leq n} c_{\alpha}(f) \varphi_{\alpha}.$$
    For $n \in \NN$, we introduce \textbf{Dirichlet kernel}
    $$\delta_{n} \coloneq \sum_{\alpha \in N, \|\alpha\| \le n} \varphi_{\alpha}.$$
\end{definitionenv}

\begin{propositionenv}\label{prop:8.5.4}
    Let $f \in L^1(\mathbb{T}_{M})$.
    One has
    $$S_n(f) = \delta_{n} * f.$$
\end{propositionenv}
\begin{proofenv}
    For any $\alpha \in N$,
    $$(\varphi_{\alpha} * f)(x) = \int_{\mathbb{T}_{M}} f(y) \varphi_{\alpha}(x - y) \eta(\dd y) = \varphi_{\alpha}(x) \int_{\mathbb{T}_{M}}f(y) \varphi_{\alpha}(-y)\eta(\dd y).$$
\end{proofenv}

\begin{remark}\label{rem:8.5.5}
    One can compute $\delta_n$ in the following way
    $$\delta_n = \sum_{(b_1,\cdots,b_d) \in \{-n,\cdots,n\}^d} \varphi_{b_1e^\vee_1 + \cdots + b_de^\vee_d} = \sum_{(b_1,\cdots,b_d) \in \{-n,\cdots,n\}^d} \prod_{j = 1}^{d} \varphi_{j}^{b_j} = \prod_{j=1}^{d} \sum_{k=-n}^{n} \varphi_j^k.$$
    Note that for any $j \in \{1,\cdots,d\}$,
    $$\sum_{k=-n}^{n} \varphi_j^k (x) = \begin{cases}
        2n+1, & x = 0\\
        \frac{\varphi_j(x)^{n+1} - \varphi_{j}(x)^{-n}}{\varphi_{j}(x) - 1}, & x \neq 0
    \end{cases}.$$
    For any $j \in \{1,\cdots,d\}$, put
    $$\delta_{n,j} = \sum_{k=-n}^{n} \varphi_j^k.$$
    By the fact that $\varphi_{j}^{-k} = \bar{\varphi_{j}^{k}}$, we obtain $\delta_{n,j}$ takes only real values.
\end{remark}

\begin{definitionenv}
    For any subset $\Theta \subseteq \{1,\cdots,d\}$, we denote by $\tau_{\Theta}$ the mapping
    $$\begin{array}{rcl}
        E & \longrightarrow & E\\
        a_1 e_1 + \cdots + a_d e_d & \longmapsto & (-1)^{\mathbbm{1}_{\Theta}(1)} a_1 e_1 + \cdots + (-1)^{\mathbbm{1}_{\Theta}(d)} a_d e_d.
    \end{array}$$
    Note that for $x = a_1 e_1 + \cdots + a_d e_d \in E$,
    $$\varphi_{j} (\tau_{\Theta}(x)) = \exp(2 \pi \mathrm{i} (-1)^{\mathbbm{1}_{\Theta}(j)} a_j) = \varphi_{j} (x)^{(-1)^{\mathbbm{1}_{\Theta}(j)}}.$$
    In particular,
    $$\delta_{n,j} (\tau_{\Theta}(x)) = \sum_{k= - n}^{n} \varphi_{j}(x)^{k(-1)^{\mathbbm{1}_{\Theta}(j)}} = \delta_{n,j}(x).$$
\end{definitionenv}

\begin{lemmaenv}[Riemann-Lebesgue]\label{lem:8.5.7}
    Let $f : \RR \longrightarrow \RR$ be a Lebesgue integrable mapping.
    Then
    $$\lim_{\lambda \to + \infty} \int_{\RR} f(t) \exp(\mathrm{i} \lambda t) \dd t = 0.$$
\end{lemmaenv}
\begin{proofenv}
    Let $a,b \in \RR$, $a<b$.
    We first consider the case where $f = \mathbbm{1}_{\interval[open left]{a}{b}}$.
    $$\int_{\RR} f(t) \exp(\mathrm{i} \lambda t) \dd t = \left[\frac{\exp(\mathrm{i \lambda t})}{\mathrm{i} \lambda}\right]^{b}_{a} = \frac{\exp(\mathrm{i} \lambda b) - \exp(\mathrm{i}\lambda a)}{\mathrm{i} \lambda}.$$
    Therefore, for $f = \mathbbm{1}_{\interval[open left]{a}{b}}$,
    $$\lim_{\lambda \to + \infty} \int_{\RR} f(t) \exp(\mathrm{i} \lambda t) \dd t = 0.$$
    By linearity, we obtain that the statement holds for any $f \in S$, where $S$ is the vector subspace of $\RR^\RR$ generated by elements of the form $\mathbbm{1}_{\interval[open left]{a}{b}}$.
    By lemma \ref{lem:7.2.4}, for any Lebesgue integrable $f$ and any $\varepsilon > 0$, there exists $g \in S$ such that
    $$\int_{\RR} |f(t) - g(t)| \dd t < \varepsilon.$$
    Hence, for any $\lambda \in \RR$,
    $$\left|\int_{\RR} f(t) \exp(\mathrm{i} \lambda t) \dd t - \int_{\RR} g(t) \exp(\mathrm{i} \lambda t) \dd t\right| \le \varepsilon.$$
    Therefore,
    $$\limsup_{\lambda \rightarrow +\infty} \left|\int_{\RR} f(t) \exp(\mathrm{i} \lambda t) \dd t\right| \le \left|\lim_{\lambda \rightarrow + \infty} \int_{\RR} g(t) \exp(\mathrm{i} \lambda t) \dd t\right| + \varepsilon.$$
\end{proofenv}

\begin{theoremenv}[Dirichlet]
    We denote by $P$ the parallelepiped
    $$\{a_1e_1 + \cdots + a_d e_d : a_1,\ldots,a_d \in \interval[open]{0}{\frac{1}{2}}\}.$$
    \begin{enumerate}[label = (\arabic*)]
        \item Let $f \in L^1(\mathbb{T}_{M})$.
            By abuse of notation, we still use $f$ to denote $f \circ \pi$ (where $\pi : E \longrightarrow \mathbb{T}_{M}$ is the projection), and thus consider $f$ as a $M$-periodic mapping on $E$.
            For any $n \in \NN$, any $x \in E$,
            $$S_{n}(f) (x) = \frac{1}{2^d \Vol(P)} \int_{P} \sum_{\Theta \subseteq \{1,\cdots,d\}} f(x + \tau_{\Theta}(y)) \delta_{n}(y) \mu(\dd y).$$
            (Again we consider $S_n(f)$ as an $M$-periodic mapping on $E$ identifying it with $S_n(f) \circ \pi$.)
        \item Suppose that $x$ is an element such that
            \begin{enumerate}
                \item for any $\Theta \subseteq \{1,\cdots,d\}$, the limit
                    $$f_{\Theta}(x) \coloneq \lim_{y \in P, y \to 0} f(x + \tau_{\Theta}(y))$$
                    exists
                \item the function
                    $$(y \in P) \longmapsto \frac{1}{e_1^\vee (y) + \cdots + e_d^\vee (y)} \sum_{\Theta \subseteq \{1,\cdots,d\}} \left(f(x + \tau_{\Theta}(y)) - f_{\Theta}(x)\right)$$
                    is integrable on $P$.
                    Then the sequence of $\left(S_{n}(f)(x)\right)_{n \in \NN}$ converges to
                    $$\frac{1}{2^d} \sum_{\Theta \subseteq \{1,\cdots,d\}} f_{\Theta}(x).$$
            \end{enumerate}
    \end{enumerate}
\end{theoremenv}
\begin{proofenv}
    \quad
    \begin{enumerate}[label = (\arabic*)]
        \item Let $D \coloneq \{a_1 e_1 + \cdots + a_d e_d \mid a_1, \cdots, a_d \in \interval[open right]{-\frac{1}{2}}{\frac{1}{2}}\}$.
            Then, since $\delta_{n}(y) = \delta_{n} (-y)$, we define
            $$S_n(f)(x) = \frac{1}{\Vol(D)} \int_{D} f(x + y)\delta_{n}(y) \mu(\dd y).$$
            By \ref{prop:8.5.4},
            $$\begin{aligned}
                S_{n}(f)(x) = & \frac{1}{\Vol(D)} \int_{D} f(x + y) \delta_{n}(y) \mu(\dd y)\\
                = & \frac{1}{2^d \Vol(P)} \sum_{\Theta \subseteq \{1,\cdots,d\}} \int_{P} f(x + \tau_{\Theta}(y)) \delta_{n}(\tau_{\Theta}(y))\mu(\dd y)\\
                = & \frac{1}{2^d \Vol(P)} \sum_{\Theta \subseteq \{1,\cdots,d\}} \int_{P} f(x + \tau_{\Theta}(y)) \delta_{n}(y)\mu(\dd y).
            \end{aligned}$$
        \item We apply (1) to the case where $f = 1$,
            $$S_n(f) (x) = 1 = \frac{1}{2^d \Vol(P)} \sum_{\Theta \subseteq \{1,\cdots,d\}} \int_{P} \delta_{n} (y) \mu(\dd y).$$
            By remark \ref{rem:8.5.5},
            $$\begin{aligned}
                \delta_{n}(y_1 e_1 + \cdots + y_d e_d) = & \prod_{j=1}^{d} \frac{\exp(\pi \mathrm{i} (2n + 1) y_{j}) - \exp(-\pi \mathrm{i} (2n + 1)y_{j})}{\exp(\pi \mathrm{i} y_{j}) - \exp(-\pi \mathrm{i} y_{j})}\\
                = & \prod_{j=1}^{d} \frac{\sin(\pi (2n + 1) y_j)}{\sin (\pi y_j)}
            \end{aligned}$$
            By our assumption,
            $$(y \in P) \longmapsto \frac{\sum_{\Theta \subseteq \{1,\cdots,d\}}(f(x + \tau_{\Theta}(y)) - f_{\Theta}(x))}{\sin(\pi e_1^\vee(y)) \cdots \sin(\pi e_d^\vee(y))}$$
            is integrable on $P$.
            (This is because $(y \in P) \longmapsto \frac{e_j^\vee (y)}{\sin(\pi e_j^\vee(y))}$ is bounded.)
            By Fubini's theorem, lemma \ref{lem:8.5.7} (Riemann-Lebesgue), and dominated convergence,
            $$\lim_{n \to +\infty} S_{n}(f) (x) = \frac{1}{2^d} \sum_{\Theta \subseteq \{1,\cdots,d\}} f_{\Theta}(x).$$
    \end{enumerate}
\end{proofenv}

\begin{remark}\label{rem:8.5.9}
    Notice in the case $d = 1$ the equality,
    $$1 = \frac{1}{2^d \Vol(P)} \sum_{\Theta \subseteq \{1,\cdots,d\}} \int_{P} \delta_{n}(y) \mu(\dd y),$$
    shows that
    $$\int_{0}^{\frac{1}{2}} \frac{\sin(\pi (2n + 1) t)}{\sin(\pi t)} \dd t = \frac{1}{2}.$$
\end{remark}

\begin{remark}
    For any $n \in \NN$, $j \in \{1,\cdots,d\}$, we define the average of $\delta_{k,j}$, $F_{n,j} : \mathbb{T}_{M} \longrightarrow \RR$, as $F_{n,j} \coloneq \frac{1}{n + 1} \sum_{k=0}^{n} \delta_{k,j}$.
    Note that
    $$F_{n,j} (0) = \frac{1}{n + 1} \sum_{k=0}^{n} (2k + 1) = n + 1$$
    and
    $$F_{n,j} = \frac{\varphi_{j}(x)^{n+2} - \varphi_{j}(x) + \varphi_{j}(x)^{-n} - \varphi_{j}(x)}{(n+1)(\varphi_{j}(x) - 1)^2}$$
    for $x \neq 0$.
    If by abuse of notation, we consider $F_{n,j}$ as a $M$-periodic function on $E$, then
    $$F_{n,j}(y_1 e_1 + \cdots + y_d e_d) = \frac{\exp(2 \pi (n + 2) \mathrm{i} y_{j}) + \exp (-2\pi n \mathrm{i} y_{j}) - 2 \exp (2\pi \mathrm{i} y_{j})}{(n+1)(\exp(2 \pi \mathrm{i} y_{j}) - 1)^2}.$$
    Whenever $y_1 e_1 + \cdots + y_d e_d$ does not belong to $M$.
    Then $n^\text{th}$ F\'ejer kernel is defined as the product of $F_{n,j}$ when $j$ varies in $\{1,\cdots,d\}$, (i.e. $F_n = \prod_{j=1}^{d} F_{n,j}$).
    Namely, as an $M$-periodic function on $E$,
    $$F_n(y_1 e_1 + \cdots + y_d e_d) = \prod_{j=1}^{d}\frac{\sin^2(\pi (n + 1)y_{j})}{(n+1)\sin^2(\pi y_{j})}.$$
    Similarly to remark \ref{rem:8.5.9},
    $$\int_{0}^{\frac{1}{2}} \frac{\sin^2(\pi(n+1)t)}{(n+1)\sin^2(\pi t)} \dd t = \frac{1}{2}.$$
\end{remark}

\begin{theoremenv}[F\'ejer]
    Let $f : \mathbb{T}_{M} \longrightarrow \CC$ be a continuous mapping.
    For any $n \in \NN$, let
    $$\tilde{S}_n (f) = \frac{1}{n + 1} \sum_{k=0}^{n} S_{k}(f).$$
    The sequence $\left(\tilde{S}_n(f)\right)_{n \in \NN}$ converges uniformly to $f$.
\end{theoremenv}
\begin{proofenv}
    By definition,
    $$\tilde{S}_{n}(f)(x) - f(x) = \frac{1}{2^d\Vol(P)} \int_{P} \sum_{\Theta \subseteq \{1,\cdots,d\}} (f(x + \tau_{\Theta}(y)) - f(x)) F_{n}(y) \mu(\dd y).$$
    For any $\delta \in \interval[open]{0}{\frac{1}{2}}$ and $\Psi \subseteq \{1,\cdots,d\}$, we let
    $$P_{\Psi}(\delta) = \{a_1 e_1 + \cdots + a_d e_d \mid (a_1,\cdots,a_d) \in \interval[open]{0}{\delta}^{\{1,\cdots,d\} \setminus \Psi} \times \interval[open right]{\delta}{\frac{1}{2}}^\Psi\},$$
    $$\varepsilon_1(\delta) = \sup_{\substack{a,b \in E\\\|a-b\| \le \delta}} |f(a) - f(b)|, \; \varepsilon_2(\delta) = \int_{\delta}^{\frac{1}{2}} \frac{\dd y}{\sin^2(\pi y)}.$$
    Without loss of generality, we may assume that $\mu \sim \Vol$ identifies with the Lebesgue measure on $\RR^d \cong E$, via $((a_1,\cdots,a_d) \longmapsto a_1 e_1 + \cdots + a_d e_d)$.
    $$\begin{aligned}
        & |\tilde{S}_{n} (f) - f(x)|\\
        \le & \int_{P_{\varnothing}(\delta)} 2^{d+1} \varepsilon_{1}(\delta) F_{n}(y) \mu(\dd y) + \sum_{\substack{\Psi \subseteq \{1,\cdots,d\}\\\Psi \neq \varnothing}} \|f\|_{\text{sup}} \left(\frac{1}{2}\right)^{d - \mathrm{card}(\Psi)} \varepsilon_2(\delta)^{\mathrm{card}(\Psi)}\\
        \le & 2^{d-1} \varepsilon_1(\delta) + \sum_{\substack{\Psi \subseteq \{1,\cdots,d\}\\\Psi \neq \varnothing}} 2^{d+1} \|f\|_{\text{sup}} \left(\frac{1}{2}\right)^{d - \mathrm{card}(\Psi)} \left(\frac{\varepsilon_2(\delta)}{n+1}\right)^{\mathrm{card}(\Psi)}.
    \end{aligned}$$
    Hence
    $$\limsup_{n \to +\infty} \|\tilde{S}_n(f) - f\|_{\text{sup}} \le 2^{d-1} \varepsilon_1(\delta).$$
    Taking the limit when $\delta \to 0^+$, we get
    $$\limsup_{n \to +\infty} \|\tilde{S}_n(f) - f\|_{\text{sup}} = 0.$$
\end{proofenv}

\begin{lemmaenv}\label{lem:8.5.12}
    Let $X$ be a locally compact Hausdorff space with a countable basis.
    Let $\mu$ be a Radon measure on $X$.
    For any $p \in \RR_{\ge 1}$, $\cc(X)$ is dense in $\mathcal{L}^p(X,\mathcal{B}(X),\mu)$ (the space of Borel-measurable real-valued functions on $X$ such that $\|f\|_{L^p} < +\infty$).
\end{lemmaenv}
\begin{proofenv}
    Let $V$ be a vector space over $\RR$ generated by $\mathbbm{1}_{A}$ for $A \subseteq X$ Borel such that $\mu(A) < +\infty$.
    Let $f : X \longrightarrow \RR_{\ge 0}$ be an element of $\mathcal{L}^p(X,\mathcal{B}(X),\mu)$.
    We can write $f$ as a pointwise limit of an increasing sequence $(f_n)_{n \in \NN}$ in $V$.
    The sequence $((f-f_n)^p)_{n \in \NN}$ which is dominated by $f^p$ converges to $0$ in $\mathcal{L}^{1}(X,\mathcal{B}(X),\mu)$.
    Hence
    $$\lim_{n \rightarrow + \infty} \|f - f_n\|_{L^p} = 0.$$
    It remains to check that any $\mathbbm{1}_{A}, A \in \mathcal{B}(X),\mu(A) < +\infty$ can be approximated in $\mathcal{L}^{p}(X,\mathcal{B}(X),\mu)$ by an element from $\cc(X)$.
    We know that for any $\varepsilon > 0$, there exists a compact set $K$ and an open set $U$ such that $K \subseteq A \subseteq U$ such that $\mu(U \setminus K) < \varepsilon$.
    By Urysohn's lemma, there exists a continuous function
    $$h : X \longrightarrow [0,1]$$
    such that $h(x) = 1$ for any $x \in K$, $\supp(h) \subseteq U$.
    Thus
    $$\int_{X}|\mathbbm{1}_{A} - h|^p \dd \mu \le \mu(U \setminus K) \le \varepsilon.$$
    Hence $\cc(X)$ is dense in $\mathcal{L}^p(X,\mathcal{B}(X),\mu)$.
\end{proofenv}

\begin{theoremenv}
    Let $f$ be an element of $\mathcal{L}^2(\mathbb{T}_{M})$.
    Then the sequence $(S_n(f))_{n \in \NN}$ converges to $f$ in $L^2(\mathbb{T}_{M})$.
    In particular, $(\varphi_{\alpha})_{\alpha \in \NN}$ forms an orthogonal basis of $L^2(\mathbb{T}_{M})$.
\end{theoremenv}
\begin{proofenv}
    First we treat the case where $f : \mathbb{T}_{M} \longrightarrow \CC$ is continuous.
    By definition, $S_{n}(f)$ is the orthogonal projection of $f$ to the vector subspace generated by $\varphi_{\alpha}$, $\|\alpha\| \le n$.
    Hence,
    $$\|f - S_{n}(f)\| \le \|f - \tilde{S}_{n}(f)\|_{L^2} \le \|f - \tilde{S}_{n}(f)\|_{\mathrm{sup}}.$$
    Passing to the limit we get
    $$\lim_{n \rightarrow +\infty} \|S_{n}(f) - f\|_{L^2} = 0.$$
    
    \quad We now consider the general case.
    By lemma \ref{lem:8.5.12}, for any $\varepsilon > 0$, there exists $g \in \cc(\mathbb{T}_{M})$ such that $\|f - g\|_{L^2} \le \varepsilon$.
    Recall that the operator norm of an orthogonal projection equals $1$.
    We deduce
    $$\|S_{n}(f) - S_{n}(g)\|_{L^2} = \|S_{n}(f - g)\|_{L^2} \le \|S_{n}\| \|f - g\|_{L^2} \le \varepsilon.$$
    Using triangle inequality,
    $$\begin{aligned}
        \|S_{n}(f) - f\|_{L^2} \le & \|S_{n}(g) - g\|_{L^2} + \|S_n(f) - S_{n}(g)\|_{L^2} + \| f - g\|_{L^2}\\
        \le & \|S_n(g) - g\|_{L^2} + 2 \varepsilon.
    \end{aligned}$$
    Hence
    $$\limsup_{n \to +\infty} \|S_{n}(f) - f\|_{L^2} \le 2 \varepsilon.$$
    Since $\varepsilon > 0$ is arbitrary,
    $$\limsup_{n \to +\infty} \|S_{n}(f) - f\|_{L^2} = 0.$$
\end{proofenv}

\begin{corollaryenv}[Parseval's equality]
    Let $f \in L^2(\mathbb{T}_{M})$.
    One has
    $$\|f\|^2_{L^2} = \sum_{\alpha \in \NN} |c_{\alpha} (f)|^2.$$
\end{corollaryenv}
\begin{proofenv}
    By definition,
    $$\|S_{n}(f)\|_{L^2}^2 = \sum_{\substack{\alpha \in \NN\\ \alpha \le n}} |c_{\alpha} (f)|^2.$$
    Take $n \to +\infty$ and we get the result.
\end{proofenv}


\section{Fourier Series on the Real Line}


In this section, we fix $E = \RR$, $M = \ZZ T$, $N = \ZZ T^{-1}$, where $T \in \RR_{> 0}$.
The results obtained in the previous section lead to the following statements.
\begin{enumerate}[label=(\arabic*)]
    \item The functions $\varphi_{n}(t) = \exp(2\pi\mathrm{i}nt/T)$, $n \in \ZZ$ form an orthogonal basis of $L^2(\RR/\ZZ T)$ equipped with the inner product induced by 
        $$\forall f,g \in \mathcal{L}^2(\RR/\ZZ T), \; \langle f,g \rangle_{L^2} = \frac{1}{T} \int_{0}^{T} f(t) \bar{g}(t) \dd t.$$
    \item For any $f \in \mathcal{L}^{2}(\RR/\ZZ T)$ and any $n \in \ZZ$, the $n$-th Fourier coefficient of $f$ is given by
        $$c_{n}(f) = \frac{1}{T} \int_{0}^{T} f(t) \exp(-2\pi\mathrm{i}n t/T) \dd t.$$
        The Fourier series of $f$ is defined as
        $$S_n(f) \coloneq \sum_{k=-n}^{n} c_{k}(f) \varphi_{k}.$$
        In the case when $f$ takes only real values,
        $$\forall n \in \ZZ, c_{-n}(f) = \overline{c_{n}(f)}.$$
        In particular, $c_{0}(f) \in \RR$ and for any positive integer $n$,
        $$\begin{aligned}
            & c_n(f) \varphi_n(t) + c_{-n}(f) \varphi_{-n}(t) = 2 \mathrm{Re}(c_{n}(f) \varphi_{n}(t))\\
            = & 2 \mathrm{Re}(c_{n}(f)) \cos(2 \pi n t/T) - 2 \mathrm{Im}(c_{n}(f)) \sin(2 \pi n t/T).
        \end{aligned}$$
        We rewrite $S_n(f)$ as
        $$S_n(f) = a_0(f) + \sum_{k=1}^{n} (a_{k}(f) \cos(2 \pi k t/T) + b_{k}(f) \sin(2 \pi k t/T)),$$
        where
        $$a_0(f) = \frac{1}{T} \int_{0}^{T} f(t) \dd t,$$
        $$a_{k}(f) = \frac{2}{T} \int_{0}^{T} f(t) \cos(2 \pi k t/T) \dd t,$$
        $$b_{k}(f) = \frac{2}{T} \int_{0}^{T} f(t) \sin(2 \pi k t/T) \dd t.$$
    \item For any $f \in \mathcal{L}^{2}(\RR/\ZZ T)$, the Fourier series of $f$
        $$S_n(f) = \frac{1}{T} \sum_{k=-n}^{n} \left(\int_{0}^{T}\exp(-2\pi\mathrm{i}kt/T) f(t) \dd t\right) \varphi_{k}$$
        converges in $L^2(\RR/\ZZ T)$ to $f$.
    \item Let $f : \RR/\ZZ T \longrightarrow \CC$ be a mapping.
        We view it as a $T$-periodic mapping from $\RR$ to $\CC$.
        We assume that
        $$f(x +) \coloneq \lim_{\substack{\varepsilon \to 0\\\varepsilon > 0}} f(x + \varepsilon),\quad f(x-) \coloneq \lim_{\substack{\varepsilon \to 0\\\varepsilon > 0}} f(x - \varepsilon)$$
        exist, and the function
        $$(y \in \interval[open]{0}{\frac{1}{2}}) \longmapsto \frac{f(x + y) + f(x - y) - f(x + ) - f(x - )}{y}$$
        is integrable with respect to Lebesgue measure.
        Then the sequence $(S_{n}(f)(x))_{n \in \NN}$ converges to $\frac{f(x+) + f(x-)}{2}$.
\end{enumerate}

\begin{propositionenv}
    Let $f : \RR /\ZZ T \longrightarrow \CC$ be a continuous mapping.
    Assume that viewed as a $T$-periodic mapping from $\RR$ to $\CC$, $f$ is of class $C^1$.
    Then
    $$c_n(f') = \frac{2 \pi \mathrm{i} n c_{n}(f)}{T}.$$
\end{propositionenv}
\begin{proofenv}
    We use integration by parts,
    $$\begin{aligned}
        c_n(f') = & \frac{1}{T} \int_{0}^{T} f'(t) \exp(-2 \pi \mathrm{i} n t /T) \dd t\\
        = & \left[\frac{f(t) \exp(-2 \pi \mathrm{i} n t/T)}{T}\right]_{0}^{T} - \frac{1}{T} \int_{0}^{T} f(t) \frac{-2 \pi \mathrm{i} n}{T} \exp(- 2 \pi \mathrm{i} n t/T) \dd t\\
        = & \frac{1}{T} \left(f(T) - f(0) + 2 \pi \mathrm{i} n c_n (f)\right)\\
        = & \frac{2 \pi \mathrm{i} n c_n(f)}{T}.
    \end{aligned}$$
\end{proofenv}

\begin{propositionenv}[Wirtinger's Inequality]
    Let $f : \RR \longrightarrow \CC$ be a $T$-periodic function of class $C^1$.
    Let
    $$m_f = \frac{1}{T} \int_{0}^{T} f(t) \dd t.$$
    Then the following inequality holds.
    $$\int_{0}^{T} |f(t) - m_f|^2 \dd t \le \frac{T^2}{4 \pi^2} \int_{0}^{T} |f'(t)|^2 \dd t.$$
\end{propositionenv}
\begin{proofenv}
    By adding a constant function to $f$, we may assume that $m_f = 0$.
    One has
    $$c_n(f') = \frac{2 \pi \mathrm{i} n c_n(f)}{T}.$$
    Notice that $c_0(f) = c_0(f') = 0$,
    $$\begin{aligned}
        \frac{1}{T} \int_{0}^{T} |f'(t)|^2 \dd t = & \sum_{n \in \ZZ \setminus \{0\}} |c_n(f')|^2 = \sum_{n \in \ZZ \setminus \{0\}} \frac{4 \pi^2 n^2}{T^2} |c_n (f)|^2\\
        \ge & \frac{4 \pi^2}{T^2} \sum_{n \in \ZZ \setminus \{0\}} |c_n(f)|^2 = \frac{4 \pi^2}{T^2} \frac{1}{T} \int_{0}^{T} |f(t)|^2 \dd t.
    \end{aligned}$$
\end{proofenv}


\section{Fourier Transform}


\begin{definitionenv}
    Let $G$ be an abelian locally compact Hausdorff group.
    A \textbf{character} of $G$ is a continuous group homomorphism
    $$\chi : G \longrightarrow S^1$$
    where $S^1 = \{z \in \CC \mid |z| = 1\}$.
    The Pontryagin dual $\hat{G}$ of $G$ is the group of characters of $G$ with operation of pointwise multiplication
    $$(\chi_1 \chi_2)(x) = \chi_1(x) \chi_2(x).$$
    Neutral element $1^\vee_{G} (x) = 1$, inverse element $\chi^{-1} (x) = \overline{\chi}(x)$.
    We give $\hat{G}$ the structure of topological group giving it so-called \textbf{compact-open} topology.
    For any pair $(K,V)$ of subsets of $G$ and $S^1$ respectively such that $K \subseteq G$ is compact and $V$ is an open neighborhood of $1$ in $S^1$, we define a subset of $\hat{G}$
    $$W(K,V) = \{\chi \in \hat{G} \mid \chi(K) \subseteq V\}.$$
    One can check that these form a neighborhood basis for $1^\vee_{G}$.
    It determines a topology on $\hat{G}$ called compact-open topology.
\end{definitionenv}

\begin{remark}
    If $G$ is discrete, the compact-open topology on $\hat{G}$ agrees with the topology of pointwise convergence.
\end{remark}

\begin{definitionenv}
    Let $G$ be an abelian locally compact Hausdorff group and let $\mu$ be a Haar measure on $G$.
    Let $f \in L^1(G,\CC)$.
    We define the \textbf{Fourier transform} $\hat{f} : \hat{G} \longrightarrow \CC$ of $f$ by the formula
    $$\hat{f}(\chi) = \int_{G} f(y) \overline{\chi}(y)\mu(\dd y).$$
\end{definitionenv}

\begin{definitionenv}
    We say that $G$ is \textbf{Pontryagin self-dual} when $G \cong \hat{G}$.
\end{definitionenv}

\begin{exampleenv}
    For $G = \RR$, consider
    $$\begin{array}{rcl}
        \RR & \longrightarrow & \hat{\RR}\\
        t & \longmapsto & (s \longmapsto \exp(2 \pi \mathrm{i} s t)).
    \end{array}$$
    It is easy to see that it is an isomorphism of groups.
    One has to show that it is a homeomorphism when we consider $\hat{\RR}$ (i.e. $\Hom_{\text{continuous}}(\RR,S^1)$ with compact-open topology).
    We can then conclude that $\RR$ is self-dual.
    The Fourier transform takes the following form
    $$\hat{f} (s \longmapsto \exp(2 \pi \mathrm{i} s t)) \sim \hat{f}(t) = \int_{\RR} f(s) \exp(-2 \pi \mathrm{i} t s) \mu(\dd s).$$
\end{exampleenv}

\begin{remark}
    We conclude that $\RR^d$ is also Pontryagin self-dual.
    Any finite-dimensional vector space over $\RR$ is Pontryagin self-dual.
\end{remark}

\begin{exampleenv}
    We fix a finite-dimensional vector space $E$ over $\RR$ and let $M$ be the sub-$\ZZ$-module of $E$ generated by a basis $(e_i)_{i=1}^{d}$ of $E$.
    Let $N$ be a sub-$\ZZ$-module of $E^\vee$ generated by $(e^\vee_i)_{i=1}^{d}$.
    $\mathbb{T}_{M} = E/M$ is considered with Haar measure $\eta$ such that $\eta(\mathbb{T}_{M}) = 1$.
    Let $f \in L^1(\mathbb{T}_{M})$, then for any $\alpha \in N$,
    $$\hat{f} (\varphi_{\alpha}) = c_{\alpha}(f) = \int_{\mathbb{T}_{M}} f(x) \overline{\varphi_{\alpha}(x)} \eta(\dd x).$$
    One can check that $N$ is Pontryagin dual of $\mathbb{T}_{M}$ (i.e. $N = \hat{\mathbb{T}}_{M}$).

    \quad In general, one can speak of Fourier series of a $M$-periodic function when we deal with a compact $E/M$.
    From the general theory, Pontryagin dual of a compact group is discrete.
\end{exampleenv}
